\documentclass[11pt]{article}

\usepackage[margin=25mm]{geometry}

\usepackage{amsmath,amssymb,fontspec,unicode-math}

\IfFontExistsTF{DejaVu Serif}{\setmainfont{DejaVu Serif}}{\setmainfont{Latin Modern Roman}}

\IfFontExistsTF{Latin Modern Math}{\setmathfont{Latin Modern Math}}{\setmathfont{DejaVu Math TeX Gyre}}

\usepackage{xurl,hyperref}

\setlength{\parindent}{0pt}\setlength{\parskip}{6pt}

\begin{document}

\title{MotionRAG: Motion Retrieval-Augmented Image-to-Video Generation}

\author{ELVIA preliminary AI analysis}\date{21 October 2026}\maketitle

Preliminary assessment: Minor.\par

Independent scientific review remains pending. Author adoption has not been requested.

The released CAMA inference branch changes the target-image range relative to training/reference frames. Its learnable initial token is also absent from the stage-\(2\) optimizer and incremental checkpoint.

\section*{Review scope}

Rapid current-version reassessment of every substantive concern in the supplied LaTeX report: the implicated paper passages and relevant author-released source paths were inspected, existing proposed corrections were briefly checked, and impact was calibrated to the demonstrated scope. No author-provided or historical verification program was executed.

\section*{Problem 1: CogVideoX inference changes the target image’s pixel range before CAMA}

Source: Section 3.3 and Figure 3, PDF p. 5; Section 3.5, PDF p. 6; released CogVideoX CAMA input path.

Assessment: supported; impact: minor

\subsection*{Source claim (paraphrase)}

Paraphrase: Section 3.3, page 5, describes the target image and the first frame of each reference video as inputs to the same DINO appearance encoder and resampler. Figure 3 depicts that common image-encoding path. Section 3.5, page 6, uses the trained CAMA component during inference. No separate target-only brightness transform is specified.

The current src/data/dataset.py, lines 88–96, scales physical pixels to \([0,1]\) and then normalizes with mean and standard deviation \(0.5\), producing \([-1,1]\). Training target frames and reference first frames follow that convention. The shared inherited DINO preprocessing maps \(x\) to \((x+1)/2\) before channel standardization.

\subsection*{Analysis}

Let a dataset pixel be \(x=2u-1\), where \(u\in[0,1]\). The common DINO preprocessing restores it by \(P(x)=(x+1)/2=u\). The inference wrapper first restores \(u\), then applies the same preprocessing again, yielding \(P(u)=(u+1)/2\). Thus black becomes \(0.5\) in the target branch but stays \(0\) in training/reference branches. A constant image makes this comparison independent of resizing.

\subsection*{Finding and scope}

A black inference target reaches the shared encoder as gray while training/reference black remains black; generated-video effects and use of that release path in published runs remain unmeasured.

\subsection*{Proposed correction}

Make the helper’s input range explicit. For the inference caller, convert \([0,1]\) back to \([-1,1]\) only before CAMA; preserve the base generator’s \([0,1]\) image. The direct training helper already receives \([-1,1]\), so an unconditional conversion in the shared helper is an invalid general patch.

The two-part local proposal remains applicable to the current release; no author adoption is established. The separate author clarification about prior short-video base adaptation is not a correction of these findings.

\section*{Problem 2: The initial context token is neither optimized nor retained by the released stage-\(2\) recipe}

Source: Figure 3, PDF p. 5; Section 3.5, PDF p. 6; Appendix A, PDF p. 21.

Assessment: supported; impact: minor

\subsection*{Source claim (paraphrase)}

Figure 3, PDF p. 5, calls the initial context a learnable token. Current ActionTransformer.\_\_init\_\_, lines 256–262, creates root parameter sos\_token with shape \((1,q,d)\), where \(q=25\), \(d=1024\), and forward line 298 uses it. The released stage-two recipe selects trainable child modules transformer and condition\_proj; the root parameter belongs to neither.

\subsection*{Analysis}

The root parameter sos\_token belongs to neither configured trainable child module transformer nor condition\_proj. Selecting only those modules omits the token from optimizer groups, and filtering checkpoint keys by the same prefixes omits it from saved state. A requires\_grad flag cannot compensate for exclusion from the optimizer and checkpoint.

\subsection*{Finding and scope}

The initial token is used but cannot be updated or reproduced from the filtered checkpoint under the current recipe; the effect of learning it on the central method is not measured.

A distributed checkpoint could contain additional state; none was downloaded. Adding a future save rule cannot reconstruct a missing historical token value.

\subsection*{Proposed correction}

The proposed correction adds the actual root parameter to optimization and serialization. A missing historical token cannot be recovered merely by changing the save filter.

A local correction or qualification is proposed where specified; author adoption is not established.

The two-part local proposal remains applicable to the current release; no author adoption is established. The separate author clarification about prior short-video base adaptation is not a correction of these findings.

\section*{Limits}

No author-provided program, training procedure, model inference or wet-lab experiment was executed in this review.\par

Historical execution results are retained only with that attribution; a source-level defect alone does not identify the program or configuration used for the published experiments.\par

No author contact, adoption, independent human scientific review or novelty determination is established.\par

The findings apply only to the identified claims, source versions and stated assumptions; they do not certify or invalidate the entire paper.\par

An existing author clarification credits prior short-video adaptation of the frozen CogVideoX base; no unequal-baseline error is inferred from the release configuration alone. Nonuniform FVD changes and the empirical-oracle interpretation remain bounded scope notes rather than newly measured errata.\par

\section*{Source identity}

Current NeurIPS 2025 final proceedings; arXiv:2509.26391v1 comparison and author code head 185f6fb5d2bd5c7ea561f1cd63dac8cd5db39af0

\url{https://proceedings.neurips.cc/paper_files/paper/2025/file/fd34c036d2408625ff7afcb4af134475-Paper-Conference.pdf}

Reviewed PDF SHA-256: \nolinkurl{b3fa4086f3aad85d3cd8457a39dcfd7b4df12f21d34f9974241498129b15f6cf}

The exact inspected current paper PDF is identified by its SHA-256. The current arXiv history lists September 2025 v1. Its implicated main-text passages match the expanded final proceedings; the formal appendix locators govern this note. The current author repository head remains the one directly inspected. An exact PDF acquisition timestamp is not established in this reassessment; source-currentness checks have their own recorded timestamps and file modification times were not substituted.

\end{document}
